\section{Diseño e Implementación del Sistema}
\label{sec:implementacion}

\subsection{Arquitectura Modular del Software}
El sistema computacional ha sido concebido bajo una arquitectura en capas estrictamente desacoplada, orientada a objetos y funciones puras, garantizando que el núcleo de cálculo de caminos mínimos sea completamente independiente de la representación visual, el almacenamiento en disco o la estructura concreta de cola de prioridad empleada (Figura \ref{fig:arquitectura}).

\begin{figure}[H]
    \centering
    \begin{tikzpicture}[
        layer/.style={rectangle, very thick, minimum width=14.8cm, minimum height=2.00cm, rounded corners=6pt},
        box/.style={rectangle, draw=black!60, fill=white, thick, minimum height=0.92cm, text width=4.30cm, align=center, rounded corners=4pt, inner ysep=4pt, inner xsep=4pt},
        boxwide/.style={rectangle, draw=black!60, fill=white, thick, minimum height=0.92cm, text width=6.65cm, align=center, rounded corners=4pt, inner ysep=4pt, inner xsep=6pt},
        arrow/.style={-{Latex[length=3.5mm, width=2.2mm]}, line width=1.3pt, draw=black!75}
    ]
    % Capa 4
    \node[layer, fill=purple!5, draw=purple!70] (L4) at (0, 8.70) {};
    \node[anchor=north, yshift=-4pt, text=purple!85!black, font=\scriptsize\bfseries] at (L4.north) {4. CAPA DE VALIDACIÓN, EXPERIMENTACIÓN Y GENERACIÓN DE REPORTES};
    \node[box] at (-4.7, 8.52) {{\fontsize{8.5}{10.5}\selectfont\textbf{\texttt{tests/test\_dijkstra.py}}}\\[2pt]{\fontsize{7.5}{9.5}\selectfont\textmd{Suite Automatizada AAA}}};
    \node[box] at (0, 8.52) {{\fontsize{8.5}{10.5}\selectfont\textbf{\texttt{scalability\_benchmark.py}}}\\[2pt]{\fontsize{7.5}{9.5}\selectfont\textmd{Banco de Escalabilidad}}};
    \node[box] at (4.7, 8.52) {{\fontsize{8.5}{10.5}\selectfont\textbf{\texttt{latex\_exporter.py}}}\\[2pt]{\fontsize{7.5}{9.5}\selectfont\textmd{Exportador de Tablas APA}}};

    % Capa 3
    \node[layer, fill=green!5, draw=green!60!black] (L3) at (0, 5.80) {};
    \node[anchor=north, yshift=-4pt, text=green!50!black, font=\scriptsize\bfseries] at (L3.north) {3. CAPA DE ALGORITMOS DE CAMINOS MÍNIMOS (SSSP)};
    \node[boxwide] at (-3.65, 5.62) {{\fontsize{8.5}{10.5}\selectfont\textbf{\texttt{src/algorithms/dijkstra.py}}}\\[2pt]{\fontsize{7.5}{9.5}\selectfont\textmd{Motor Canónico de Relajación Voraz}}};
    \node[boxwide] at (3.65, 5.62) {{\fontsize{8.5}{10.5}\selectfont\textbf{\texttt{dijkstra\_solver(graph, \dots)}}}\\[2pt]{\fontsize{7.5}{9.5}\selectfont\textmd{Inyección Abstracta de Colas de Prioridad}}};

    % Capa 2
    \node[layer, fill=orange!6, draw=orange!75!black] (L2) at (0, 2.90) {};
    \node[anchor=north, yshift=-4pt, text=orange!85!black, font=\scriptsize\bfseries] at (L2.north) {2. CAPA DE ESTRUCTURAS DE DATOS Y COLAS DE PRIORIDAD};
    \node[box] at (-4.7, 2.72) {{\fontsize{8.5}{10.5}\selectfont\textbf{\texttt{IndexedMinHeap}}}\\[2pt]{\fontsize{7.5}{9.5}\selectfont\textmd{Min-Heap Binario Indexado}}};
    \node[box] at (0, 2.72) {{\fontsize{8.5}{10.5}\selectfont\textbf{\texttt{DAryHeap} ($d=4$)}}\\[2pt]{\fontsize{7.5}{9.5}\selectfont\textmd{Montículo 4-ario en Caché}}};
    \node[box] at (4.7, 2.72) {{\fontsize{8.5}{10.5}\selectfont\textbf{\texttt{FibonacciHeap}}}\\[2pt]{\fontsize{7.5}{9.5}\selectfont\textmd{Cortes en Cascada $O(1)^*$}}};

    % Capa 1
    \node[layer, fill=blue!5, draw=blue!70!black] (L1) at (0, 0.0) {};
    \node[anchor=north, yshift=-4pt, text=blue!85!black, font=\scriptsize\bfseries] at (L1.north) {1. CAPA DE MODELOS TOPOGRÁFICOS, COSTOS Y SERIALIZACIÓN};
    \node[box] at (-4.7, -0.18) {{\fontsize{8.5}{10.5}\selectfont\textbf{\texttt{cusco\_network.py}}}\\[2pt]{\fontsize{7.5}{9.5}\selectfont\textmd{Topología Vial $G=(V,E)$}}};
    \node[box] at (0, -0.18) {{\fontsize{8.5}{10.5}\selectfont\textbf{\texttt{cost\_model.py}}}\\[2pt]{\fontsize{7.5}{9.5}\selectfont\textmd{Impedancia Andina $\Phi(\theta)$}}};
    \node[box, fill=gray!10] at (4.7, -0.18) {{\fontsize{7.8}{9.5}\selectfont\textbf{\texttt{cusco\_centro\_network.json}}}\\[2pt]{\fontsize{7.5}{9.5}\selectfont\textmd{Red Vial Serializada (6 nodos)}}};

    % Flechas con espacio holgado y etiquetas limpias
    \draw[arrow] (L4.south) -- (L3.north) node[midway, right=8pt, font=\scriptsize\itshape, text=black!75] {ejecuta e instrumenta};
    \draw[arrow] (L3.south) -- (L2.north) node[midway, right=8pt, font=\scriptsize\itshape, text=black!75] {invoca operaciones de cola};
    \draw[arrow] (L2.south) -- (L1.north) node[midway, right=8pt, font=\scriptsize\itshape, text=black!75] {consume pesos e incidencias};
    \end{tikzpicture}
    \caption{Diagrama arquitectural desacoplado del sistema en cuatro capas operativas}
    \label{fig:arquitectura}
\end{figure}

\subsection{Descripción de Módulos del Sistema}
El repositorio organiza su código fuente en el paquete \texttt{src/}:
\begin{enumerate}
    \item \textbf{\texttt{src/algorithms/dijkstra.py}:} Implementa el solver de caminos mínimos con inyección de dependencias. Recibe la lista de adyacencia y una cadena identificadora del tipo de cola (\texttt{binary\_heap}, \texttt{4ary\_heap}, \texttt{fibonacci\_heap}), ejecutando las relajaciones y capturando las estadísticas de operaciones para su posterior análisis.
    \item \textbf{\texttt{src/datastructures/indexed\_min\_heap.py}:} Implementación pura del Min-Heap binario con tabla de posiciones invertida para garantizar \texttt{decrease\_key} en tiempo $O(\log |V|)$.
    \item \textbf{\texttt{src/datastructures/d\_ary\_heap.py}:} Implementación del montículo $d$-ario parametrizado para cualquier aridad $d \ge 2$, configurado por defecto en $d=4$.
    \item \textbf{\texttt{src/datastructures/fibonacci\_heap.py}:} Implementación del Fibonacci Heap con nodos enlazados circularmente, cálculo de rangos y corte en cascada.
    \item \textbf{\texttt{src/models/cusco\_network.py}:} Mapeo explícito de la topología vial del Centro Histórico de Cusco, serialización y carga en formato JSON.
    \item \textbf{\texttt{src/utils/cost\_model.py}:} Lógica de cálculo del costo generalizado e impedancia por gradiente altitudinal andino.
    \item \textbf{\texttt{src/utils/latex\_exporter.py}:} Rutina automatizada que formatea diccionarios de traza en tablas LaTeX bajo el estándar \texttt{booktabs}.
    \item \textbf{\texttt{src/visualization/generate\_figures.py}:} Renderizador de gráficos vectoriales y matriciales de alta resolución (300 DPI) con Matplotlib y NetworkX.
\end{enumerate}

\subsection{Formatos de Datos y Serialización JSON}
Para garantizar que la ingestión de datos no dependa de estructuras en memoria volátiles ni rutas locales rígidas, la red vial del Cusco se serializa formalmente en formato JSON bajo la ruta estándar \texttt{data/processed/cusco\_centro\_network.json}:
\begin{lstlisting}[language=Python, caption={Extracto del esquema JSON para la red vial del Centro Histórico de Cusco}]
{
  "nodes": [
    {"id": "A", "name": "Plaza Mayor del Cusco"},
    {"id": "B", "name": "Plaza Regocijo (Cusipata)"}
  ],
  "edges": [
    {
      "origin": "A",
      "destination": "B",
      "length_m": 180.0,
      "street_type": "calle_centro",
      "gradient_pct": 1.5,
      "street_name": "Calle Espaderos / Portal Carnicerias",
      "travel_time_sec": 44.75
    }
  ]
}
\end{lstlisting}

\subsection{Contratos de Interfaz de las Estructuras de Datos}
Para posibilitar la sustitución dinámica y transparente de colas en el solver de Dijkstra sin alterar el flujo algorítmico, cada una de las tres estructuras implementa una interfaz uniforme:
\begin{itemize}
    \item \texttt{insert(item: str, key: float) -> None}: Inserta un identificador de vértice con su distancia tentativa asociada.
    \item \texttt{extract\_min() -> Tuple[str, float]}: Extrae y retorna el elemento con menor clave acumulada, reestructurando el montículo según sus invariantes.
    \item \texttt{decrease\_key(item: str, new\_key: float) -> None}: Disminuye la clave del vértice a su nuevo valor relajado, flotando el nodo o ejecutando cortes en cascada.
    \item \texttt{is\_empty() -> bool}: Determina en $O(1)$ si existen vértices pendientes de procesamiento.
\end{itemize}

\subsection{Implementación de Algoritmos y Estructuras Nucleares}
A continuación, se presentan los códigos fuente de los algoritmos y estructuras nucleares desarrollados, embebidos directamente desde el repositorio:

\lstinputlisting[
    language=Python,
    caption={Implementación desacoplada del algoritmo de Dijkstra con inyección de colas (\texttt{python/src/algorithms/dijkstra.py})},
    label={lst:dijkstra_py},
    firstline=1,
    lastline=75
]{../python/src/algorithms/dijkstra.py}

\lstinputlisting[
    language=Python,
    caption={Implementación del Min-Heap Binario Indexado (\texttt{python/src/datastructures/indexed\_min\_heap.py})},
    label={lst:indexed_heap_py},
    firstline=1,
    lastline=60
]{../python/src/datastructures/indexed_min_heap.py}

% [PENDIENTE RAMA YENI (Yeni Porroa - 120893): Descomentar siguiente listing tras copiar d_ary_heap.py]
\lstinputlisting[
    language=Python,
    caption={Implementación del Montículo 4-ario contiguo (\texttt{python/src/datastructures/d\_ary\_heap.py})},
    label={lst:dary_heap_py},
    firstline=1,
    lastline=60
]{../python/src/datastructures/d_ary_heap.py}

% [PENDIENTE COMMIT 3 (Romario Quispe - 164257): Descomentar siguiente listing tras copiar fibonacci_heap.py]
% \lstinputlisting[
%     language=Python,
%     caption={Implementación del Fibonacci Heap con cortes en cascada (\texttt{python/src/datastructures/fibonacci\_heap.py})},
%     label={lst:fibonacci_heap_py},
%     firstline=1,
%     lastline=70
% ]{../python/src/datastructures/fibonacci_heap.py}

\lstinputlisting[
    language=Python,
    caption={Modelo de costo e impedancia topográfica andina (\texttt{python/src/utils/cost\_model.py})},
    label={lst:cost_model_py},
    firstline=1,
    lastline=44
]{../python/src/utils/cost_model.py}
